\section{Generación de reportes}
\label{sec:marco_reportes}

\subsection{El reporte clínico dermatológico}
\label{sec:marco_rep_tarea}

Un reporte clínico dermatológico es el registro escrito del examen de una lesión. Lo redacta el médico que la
examina, queda en la historia clínica y comunica lo observado a otros profesionales, por ejemplo al
especialista al que se deriva el caso. Su contenido se organiza en tres partes
(Tabla~\ref{tab:marco_partes_reporte}): la descripción de los hallazgos visibles, la impresión diagnóstica que
los resume y, cuando corresponde, una recomendación de conducta. La descripción usa un vocabulario
estandarizado por consenso, con términos para el color, el patrón, la distribución y las estructuras de la
lesión, y algoritmos como la lista de 7 puntos convierten la presencia de ciertos hallazgos en peso
diagnóstico~\cite{kittler2016standardization,argenziano1998checklist}.

\begin{table}[H]
    \centering
    \caption{Partes de un reporte dermatológico y lo que se le exige a cada una.}
    \label{tab:marco_partes_reporte}
    \small
    \begin{tabularx}{\textwidth}{@{}>{\raggedright\arraybackslash}p{3.0cm}
        >{\raggedright\arraybackslash\hspace{0pt}}X >{\raggedright\arraybackslash\hspace{0pt}}X@{}}
        \toprule
        \textbf{Parte} & \textbf{Contenido} & \textbf{Qué se le exige} \\
        \midrule
        Descripción & Color, patrón, bordes, estructuras y localización de la lesión
            & Que describa lo que la imagen muestra, con el vocabulario de la especialidad \\
        \filasep
        Impresión diagnóstica & La categoría diagnóstica, o los diferenciales si hay duda
            & Que sea la correcta y que su grado de certeza corresponda a la evidencia \\
        \filasep
        Recomendación & Seguimiento, biopsia o derivación
            & Que sea coherente con la impresión diagnóstica \\
        \bottomrule
    \end{tabularx}
\end{table}

Generar un reporte de forma automática tiene una historia corta. Su antecedente es la descripción automática
de imágenes (\emph{image captioning}), que produce una oración sobre una imagen natural. El reporte médico es
más largo, sigue una estructura y usa un vocabulario técnico. Su generación con redes neuronales se
desarrolló primero en radiología, desde 2018, con modelos que predicen etiquetas de hallazgos y redactan el reporte oración por
oración, y con arquitecturas que recuerdan patrones de redacción entre reportes
parecidos~\cite{jing2018automatic,chen2020generating}. Con los \glspl{vlm} generativos, la tarea pasó a
resolverse con un único modelo que recibe la imagen y una instrucción~\cite{hartsock2024vision}. En
todos los casos el entrenamiento requiere pares de imagen y texto clínico, y en dermatología esos pares son
escasos y recientes: su texto fue redactado por dermatólogos, tomado de libros de la especialidad o generado
por un modelo y revisado por especialistas~\cite{shen2025skincare,zeng2025mmskin,yan2025derm1miccv}, y ese origen limita cuánto se puede
confiar en él como referencia.

\subsection{Tres ejes de la generación de reportes}
\label{sec:marco_rep_ejes}

Generar un reporte resuelve tres problemas a la vez, uno por cada parte importante del texto.

\looseness=-1 El primero es describir lo que se ve: el texto tiene que reflejar los rasgos visuales de la
lesión con el vocabulario de la especialidad. La \textbf{calidad descriptiva} es la correspondencia entre lo
que el texto dice y lo que la imagen muestra, y el vocabulario cerrado permite medirla
(Sección~\ref{sec:marco_eval_texto}).

El segundo es incorporar información diagnóstica, porque el texto tiene que nombrar una categoría. La
\textbf{fidelidad diagnóstica} es que el diagnóstico que nombra el texto sea el correcto, y es independiente de
la calidad descriptiva: un texto puede describir fielmente la imagen y cerrar en un diagnóstico equivocado, y
otro puede nombrar el diagnóstico correcto acompañado de hallazgos inventados. Un reporte es correcto solo si
cumple las dos condiciones (Figura~\ref{fig:intro_dos_tareas}).

% Las dos condiciones de un reporte clínico correcto (clasificar y describir una misma imagen).
% Se trajo desde el Capítulo 1 a sec:marco_rep_ejes (2.6.2). El nombre del archivo y la etiqueta
% (fig:intro_dos_tareas) se conservan.

\begin{figure}[H]
    \centering
    \begin{tikzpicture}[x=1mm, y=1mm]

        \node[figentrada, text width=27mm] (img) at (12, 0) {Imagen\\dermatoscópica};

        \node[figtarea, text width=36mm] (t1) at (58, 14)
            {\textbf{Clasificar}\\asignar una categoría};
        \node[figtarea, text width=36mm] (t2) at (58, -14)
            {\textbf{Describir}\\redactar el texto\\de la lesión};

        \node[figsalida, text width=34mm] (cat) at (114, 14)
            {Categoría diagnóstica\\[2pt]{\footnotesize se evalúa con la\\fidelidad diagnóstica}};
        \node[figsalida, text width=34mm] (txt) at (114, -14)
            {Descripción de la lesión\\[2pt]{\footnotesize se evalúa con la\\calidad descriptiva}};

        \draw[figflecha] (img.east) -- ++(6, 0) |- (t1.west);
        \draw[figflecha] (img.east) -- ++(6, 0) |- (t2.west);
        \draw[figflecha] (t1.east) -- (cat.west);
        \draw[figflecha] (t2.east) -- (txt.west);

        \begin{scope}[on background layer]
            \node[figgrupo={las dos condiciones de un reporte correcto},
                  fit=(t1) (t2)] {};
        \end{scope}

    \end{tikzpicture}
    \caption{Los dos requisitos de un reporte clínico correcto: la categoría diagnóstica que nombra se evalúa por su fidelidad diagnóstica, y la descripción de la lesión, por su calidad descriptiva.}
    \label{fig:intro_dos_tareas}
\end{figure}


El tercero es tratar las predicciones inciertas. Cuando el diagnóstico no está claro, afirmar uno arriesga un
error que el lector no puede detectar, y la política de la Sección~\ref{sec:marco_tres_salidas} decide si el
texto afirma, enumera diferenciales o no se redacta. Los tres problemas dependen de cómo se produce el texto.

\subsection{Cómo se genera un reporte}
\label{sec:marco_rep_generacion}

Hay dos formas de generar un reporte con un \gls{vlm} (Figura~\ref{fig:marco_generacion_reportes}). En la
generación \textbf{de extremo a extremo}, el modelo recibe la imagen $x$ y una instrucción, y produce el texto
completo, $y \sim p_\theta(\cdot \mid x)$: describe la lesión e infiere el diagnóstico a partir de lo que ve.
Resuelve los tres problemas con un solo modelo, que tiene que ser a la vez un buen clasificador y un buen
redactor.

% Figura: dos formas de generar un reporte con un VLM, de extremo a extremo o condicionado en evidencia
% externa. Se usa en sec:marco_rep_generacion.

\begin{figure}[H]
    \centering
    \begin{tikzpicture}[x=1mm, y=1mm]

        % (a) De extremo a extremo.
        \node[font=\small\bfseries, anchor=west] at (-6, 30) {(a) De extremo a extremo};
        \node[figentrada, text width=22mm] (a_img) at (8, 18) {Imagen e\\instrucción};
        \node[figtarea, text width=26mm] (a_vlm) at (62, 18) {\gls{vlm}};
        \node[figsalida, text width=26mm] (a_rep) at (122, 18) {Reporte};
        \draw[figflecha] (a_img) -- (a_vlm);
        \draw[figflecha] (a_vlm) -- (a_rep);

        % (b) Condicionado en evidencia externa. Arriba, el componente que aporta la evidencia y la regla de
        % decisión; abajo, el mismo recorrido de (a) con la evidencia entrando al VLM desde arriba.
        \node[font=\small\bfseries, anchor=west] at (-6, 4) {(b) Condicionado en evidencia externa};
        \node[figtarea, text width=28mm] (b_ev) at (38, -8) {Clasificador o\\recuperación};
        \node[figtarea, text width=34mm] (b_pol) at (84, -8)
            {Regla de decisión\\[1pt]{\footnotesize afirmar, diferencial\\o derivar}};
        \node[figentrada, text width=22mm] (b_img) at (8, -34) {Imagen};
        \node[figtarea, text width=26mm] (b_vlm) at (84, -34) {\gls{vlm}};
        \node[figsalida, text width=26mm] (b_rep) at (122, -34) {Reporte};
        \draw[figflecha] (b_img.north) |- (b_ev.west);
        \draw[figflecha] (b_ev) -- (b_pol);
        \draw[figflecha] (b_pol) -- node[figetiqueta, right=1mm] {evidencia $c(x)$} (b_vlm);
        \draw[figflecha] (b_img) -- (b_vlm);
        \draw[figflecha] (b_vlm) -- (b_rep);

    \end{tikzpicture}
    \caption{Dos formas de generar un reporte: el \gls{vlm} solo o con la evidencia de
    otro componente.}
    \label{fig:marco_generacion_reportes}
\end{figure}


En la generación \textbf{condicionada en evidencia externa}, otro componente aporta información que se agrega al
\emph{prompt} del generador,
\begin{equation}
    y \sim p_\theta\bigl(\,\cdot \mid x,\ c(x)\bigr),
    \label{eq:marco_rep_condicionada}
\end{equation}
donde $c(x)$ es esa evidencia serializada como texto. Puede ser el diagnóstico que predice un clasificador, los
diagnósticos de casos parecidos recuperados de una base (Sección~\ref{sec:marco_rep_recuperacion}) o los
candidatos que deja la política de la Ecuación~\ref{eq:marco_politica}. El generador conserva la imagen y
redacta la descripción, pero el diagnóstico le llega como dato. Anclar así el texto en evidencia externa se
conoce como \emph{grounding}, y es aprendizaje en contexto (Sección~\ref{sec:marco_llm_grandes}): el
generador no cambia sus pesos~\cite{brown2020language}.

La segunda forma separa los problemas: la fidelidad diagnóstica pasa a depender de la evidencia, y la calidad
descriptiva sigue a cargo del generador. Tiene dos costos. El reporte hereda los errores de la evidencia,
porque un generador que recibe un diagnóstico equivocado lo afirma con la misma fluidez que uno correcto, y el
sistema usa más de un modelo. La instrucción también forma parte de la configuración: puede pedir que el
modelo razone antes de concluir o que entregue solo el reporte, y con eso cambia lo que el texto afirma. En las
dos formas el texto puede fallar de maneras que una métrica de clasificación no ve.

\subsection{Fallas del texto generado}
\label{sec:marco_rep_fallas}

Una alucinación es un texto fluido cuyo contenido no está sustentado por la entrada, que en este caso es la
imagen. Como falla, el fenómeno se describió primero en la generación de lenguaje
natural~\cite{ji2023survey} y en la descripción de imágenes, donde los modelos nombran objetos que la imagen no contiene en parte por el peso de
los patrones del lenguaje~\cite{rohrbach2018object}, y se documentó después en los
\glspl{vlm}~\cite{bai2024hallucination}. En los \glspl{vlm} médicos, un estudio sistemático encontró que los
modelos ajustados al dominio son más propensos a alucinar que los de propósito general, aun cuando rinden
mejor en las tareas médicas estándar~\cite{gu2024medvh}. La alucinación es especialmente costosa en medicina,
porque el modelo describe con la misma fluidez los hallazgos que la imagen contiene y los que no, y un
hallazgo inventado puede cambiar una conducta clínica.

La falla contraria se llama aquí \emph{hedging}. El modelo redacta descripciones genéricas, compatibles con
casi cualquier lesión, y evita comprometer un diagnóstico, por ejemplo enumerando varias categorías posibles
sin inclinarse por ninguna. Ese texto no contiene errores visibles y tampoco informa nada sobre la lesión
concreta. Expresar incertidumbre no es en sí una falla: se convierte en \emph{hedging} cuando es
indiscriminada, es decir, cuando no depende de cuán incierto sea el caso. Se ha argumentado que los esquemas
de evaluación que premian acertar y no penalizan adivinar incentivan la alucinación en lugar de la expresión
de incertidumbre~\cite{kalai2025hallucinate}. Enumerar diferenciales cuando una señal de confianza lo
justifica (Sección~\ref{sec:marco_tres_salidas}) es lo opuesto al \emph{hedging}.

\begin{table}[H]
    \centering
    \caption{Fallas de un reporte generado y su consecuencia para el lector, de mayor a menor gravedad.}
    \label{tab:marco_fallas}
    \small
    \begin{tabularx}{\textwidth}{@{}>{\raggedright\arraybackslash}p{3.0cm}
        >{\raggedright\arraybackslash\hspace{0pt}}X >{\raggedright\arraybackslash\hspace{0pt}}X@{}}
        \toprule
        \textbf{Falla} & \textbf{Qué hace el texto} & \textbf{Consecuencia para el lector} \\
        \midrule
        Diagnóstico incorrecto & Afirma una categoría distinta de la real
            & Recibe un error que la fluidez del texto hace parecer confiable \\
        \filasep
        Alucinación & Describe hallazgos que la imagen no contiene
            & Puede cambiar una conducta y no se distingue por la forma del texto \\
        \filasep
        Contradicción & Nombra dos o más diagnósticos incompatibles como conclusión
            & No sabe cuál de ellos sostiene el reporte \\
        \filasep
        \emph{Hedging} & Redacta un texto genérico y evita comprometer un diagnóstico
            & No recibe información y no sabe si el modelo duda o no \\
        \filasep
        Repetición & Repite frases en bucle~\cite{holtzman2020curious}
            & Recibe un texto degenerado, sin información útil \\
        \filasep
        Texto insuficiente & Una o dos frases sin descripción de la lesión
            & No recibe la descripción que el reporte debe contener \\
        \bottomrule
    \end{tabularx}
\end{table}

Estas dos fallas no son las únicas. Al leer los reportes uno por uno, cada uno se puede asignar a su falla
principal (Tabla~\ref{tab:marco_fallas}). Afirmar un diagnóstico incorrecto es la más grave, porque el lector
recibe un error con la forma de un acierto. La contradicción, que nombra dos diagnósticos incompatibles en el
mismo texto, y el \emph{hedging} dejan la decisión en el lector, y la repetición y el texto insuficiente no
aportan información. La alucinación puede acompañar a cualquiera de ellas, incluso a un diagnóstico correcto.
La taxonomía organiza el análisis de los reportes generados del Capítulo~\ref{cp:res_gen}.

Las métricas de solapamiento con un texto de referencia (Sección~\ref{sec:marco_eval_texto}) miden un
parecido superficial y pueden no detectar si el reporte transmite los hallazgos clínicos
correctos~\cite{hartsock2024vision}, por lo que la fidelidad
diagnóstica y el contenido clínico se miden por separado. Una manera de reducir estas fallas es darle al
modelo, por la vía de la Ecuación~\ref{eq:marco_rep_condicionada}, información que no obtiene de la imagen
sola, y la recuperación es la forma más estudiada de obtenerla.

\subsection{La recuperación como fuente de evidencia}
\label{sec:marco_rep_recuperacion}

\phantomsection\label{sec:marco_rag}%
La \gls{rag}, de 2020, agrega al generador un paso de búsqueda en una base externa de documentos. Tiene tres
partes: una base de documentos, un recuperador que selecciona los más relevantes para cada entrada y un
generador que produce la salida condicionado en la entrada y en lo recuperado. En la formulación original, en su variante
RAG-Sequence, un codificador proyecta cada documento $z$ de la base a un embedding $\mathbf{d}(z)$ y otro proyecta la consulta
$x$ a $\mathbf{q}(x)$. El recuperador asigna a cada documento una probabilidad
$p_\eta(z \mid x) \propto \exp\!\left(\mathbf{d}(z)^{\top}\mathbf{q}(x)\right)$, y el generador produce la
salida $y$ con cada uno de los $K$ documentos más probables y combina los resultados pesándolos por esa
probabilidad,
\begin{equation}
    p(y \mid x) \approx \sum_{z \in \text{top-}K} p_\eta(z \mid x)\, p_\theta(y \mid x, z).
    \label{eq:marco_rag}
\end{equation}
El codificador de la consulta se ajusta junto con el generador, mientras que el de documentos y el índice
quedan fijos, y la base se puede cambiar o ampliar sin reentrenar~\cite{lewis2020retrievalaugmented}. En
radiología, la recuperación multimodal parte de la imagen y recupera reportes de casos parecidos, y entrenar el
recuperador para que prefiera reportes con los mismos hallazgos clínicos mejora la corrección de lo
generado~\cite{sun2025factaware}.

\phantomsection\label{sec:marco_recuperacion_imagenes}%
Con una base de casos etiquetados, lo recuperable son imágenes. Los $k$ vecinos de la
Ecuación~\ref{eq:marco_vecinos} son los casos más parecidos a la consulta en el espacio de un codificador
visual, y su información, como su diagnóstico o su localización, es \textbf{evidencia visual}: dice qué se vio
que se parezca a la consulta. Aplicada a un modelo de lenguaje, la búsqueda de vecinos mejora la predicción sin
entrenamiento adicional, sobre todo en los patrones poco frecuentes, porque el modelo accede a casos concretos
que no tiene que memorizar~\cite{khandelwal2020generalization}. Un antecedente en descripción de imágenes
recupera las descripciones de imágenes parecidas y las entrega como \emph{prompt} a un modelo de lenguaje, sin
entrenar el recuperador~\cite{ramos2023smallcap}.

Esta segunda forma es un caso de la Ecuación~\ref{eq:marco_rep_condicionada} con
$c(x) = \psi\bigl(\mathcal{N}_k(x)\bigr)$, donde $\psi$ serializa como texto la información de los vecinos, y
difiere de la formulación de \textcite{lewis2020retrievalaugmented} en tres aspectos.
\begin{enumerate}
    \item El recuperador puede quedar congelado, sin entrenarse junto con el generador.
    \item No se marginaliza: el generador recibe todo el contexto a la vez y produce una única salida.
    \item Lo recuperado son etiquetas y atributos de imágenes, de la misma modalidad que la consulta, y no
    documentos de texto.
\end{enumerate}

\phantomsection\label{sec:marco_integracion}%
Recuperar documentos, recuperar imágenes y medir la concordancia entre predicciones de la
Sección~\ref{sec:marco_senales_confianza} no son equivalentes (Tabla~\ref{tab:marco_mecanismos}). Los dos
primeros aportan evidencia al texto y atienden el segundo eje de la Sección~\ref{sec:marco_rep_ejes}; la
concordancia no aporta evidencia, pero decide cuánta afirmación permite la que hay y atiende el tercero.
Integrados, la evidencia recuperada condiciona el texto y la señal de confianza elige si el texto afirma,
enumera diferenciales o no se redacta.

\begin{table}[H]
    \centering
    \caption{Mecanismos que agregan información al generador: qué aportan, su modalidad, qué ayudan a resolver
    y qué dejan abierto.}
    \label{tab:marco_mecanismos}
    \small
    \begin{tabularx}{\textwidth}{@{}>{\raggedright\arraybackslash}p{2.7cm}
        >{\raggedright\arraybackslash\hspace{0pt}}X >{\raggedright\arraybackslash\hspace{0pt}}X
        >{\raggedright\arraybackslash\hspace{0pt}}X >{\raggedright\arraybackslash\hspace{0pt}}X@{}}
        \toprule
        \textbf{Mecanismo} & \textbf{Qué recupera o calcula} & \textbf{Modalidad} & \textbf{Ayuda a resolver} & \textbf{No resuelve} \\
        \midrule
        Recuperar documentos (\gls{rag}) & Textos de una base, seleccionados por un recuperador
            & Consulta y documentos; los documentos son texto
            & Descripción y diagnóstico con conocimiento textual
            & Requiere una base de documentos pertinentes y no dice cuánto confiar en ellos \\
        \filasep
        Recuperar imágenes similares & Etiquetas y atributos de los $k$ casos más cercanos
            & Imagen contra imagen
            & Diagnóstico, como evidencia visual
            & No aporta redacción, y los vecinos existen aunque no se parezcan a la consulta \\
        \filasep
        Concordancia entre predicciones & Si dos predictores coinciden sobre la misma imagen
            & Predicciones
            & Incertidumbre, es decir, cuánta afirmación se permite
            & No agrega información al texto ni corrige un error que los dos predictores comparten \\
        \bottomrule
    \end{tabularx}
\end{table}

Clasificadores, generadores y señales de confianza producen así salidas de tres tipos: categorías, textos y
decisiones de abstención. Cada una exige sus propias métricas.
